\documentclass[11pt,a4paper]{article}
\usepackage[T1]{fontenc}
\usepackage[utf8]{inputenc}
\usepackage[french]{babel}
\usepackage{lmodern}
\usepackage{amsmath,amssymb}
\usepackage[margin=2.35cm]{geometry}
\usepackage{booktabs,longtable,tabularx}
\usepackage{enumitem}
\usepackage{graphicx}
\usepackage{xcolor}
\usepackage{microtype}
\usepackage{hyperref}

\hypersetup{colorlinks=true,linkcolor=blue!45!black,urlcolor=blue!45!black}
\setlist{nosep}
\newcommand{\code}[1]{\texttt{#1}}
\newcommand{\Kaut}{K^{*}_{\mathrm{aut}}}
\newcommand{\fait}{\textbf{[Fait observé]}}
\newcommand{\inference}{\textbf{[Inférence]}}
\newcommand{\hyp}{\textbf{[Hypothèse]}}
\newcommand{\incertitude}{\textbf{[Incertitude]}}

\title{M4.2B --- rapport de phase pilote\\
\large Émergence et contrôle d'une queue Pareto des revenus d'intérêt,
sous contrainte de cascades critiques}
\author{Programme M4.2B --- dossier \code{m4\_2b\_credit\_soc}}
\date{30 juillet 2026}

\begin{document}
\maketitle

\begin{abstract}
M4.2B remplace la cible de principal géométrique de M4.2,
$K^*=\sqrt{K_\ell K_b}$, par une cible \textbf{arithmétique}
$K_{\mathrm{target}}=(K_\ell+K_b)/2$, à taux inchangé, dans une baseline
scientifique volontairement lente ($\delta=0{,}01$, $\sigma=0{,}01$). La
question posée était double : une queue Pareto endogène des intérêts reçus
peut-elle émerger et être contrôlée, et la structure en loi de puissance
des avalanches de faillites survit-elle ? Ce rapport documente une
\textbf{phase pilote}, pas une campagne figée : les runs exécutés (baseline
$T=3000$, contrôle géométrique, balayages $\sigma$ et $\lambda$,
contraste $\rho$) ont été choisis pour maximiser l'information causale
avant tout engagement de calcul dans une grille exhaustive, conformément
à la liberté de recherche accordée par le cahier des charges.

Le résultat central est \textbf{négatif pour l'objectif A, avec un
mécanisme identifié et quantifié} : à la baseline, le capital devient
presque parfaitement homogène (Gini$(K)\approx0{,}05$) --- fait qui
persiste sous les deux règles de principal et sur un facteur 10 en
population --- et le revenu d'intérêt reçu est dominé à 77\,\% par le
nombre de contrats accumulés avec l'âge, non par leur taille ou leur taux
individuels. Aucune configuration testée (baseline, $\sigma\in[0{,}01\,;\,
0{,}10]$, $\rho\in\{0{,}125\,;\,1\,;\,2\}$, $\lambda\in\{10\,;\,30\,;\,100\}$,
règle géométrique) ne produit une queue Pareto stable au sens du seuil :
l'exposant ajusté dérive systématiquement d'un facteur $\sim\!1{,}5$ quand
le seuil est divisé par deux. Fait notable et contre-intuitif : la règle
arithmétique, changement constitutif de M4.2B, rend l'objectif A
\emph{plus difficile} que la règle géométrique de M4.2 (queue plus légère,
dominance du degré plus forte), pas plus facile. Pour l'objectif B, le
rapport de branchement des avalanches est robuste et institutionnellement
pinné ($\approx0{,}78$--$0{,}79$) sur toute la plage testée, mais
\textbf{aucun exposant de queue indépendant de la taille du système n'a pu
être établi} : l'estimateur de coupure sature à sa borne numérique à
$\lambda=100$, signe que l'exposant rapporté y est en réalité celui de la
loi pure, non comparable aux cellules $\lambda=10$ et $\lambda=30$.
\end{abstract}

\tableofcontents

\section{Modification constitutive et vérifications}

\subsection{La cible arithmétique}

Pour $K_\ell>K_b>0$, M4.2B calcule directement
\[
q_A=\frac{K_\ell-K_b}{2},\qquad K_\ell'=K_b'=\frac{K_\ell+K_b}{2},
\]
sans passer par la fonction $K^*(r)$ de M4.2. Le taux reste
$r_{\ell b}=\sqrt{m_\ell m_b}$, $m=A\gamma K^{\gamma-1}$, strictement
séparé du calcul du principal dans le code (\code{\_pair\_rate} /
\code{\_pair\_principal}).

\fait\ Le rapport entre les deux institutions à taux identique est exact :
\[
\frac{q_A}{q_{\mathrm{geo}}}=\frac{1+\sqrt{K_\ell/K_b}}{2},
\]
vérifié analytiquement et sur grille $(K_\ell,K_b,\gamma,A)$ à $10^{-9}$
près (\code{tests/test\_arithmetic\_institution.py}). Ce rapport n'est
borné par rien : pour $K_\ell/K_b=392$ (un newborn $K_0=25$ face à une
entité proche de $\Kaut(\gamma{=}0{,}5)=9801$), il vaut $10{,}4$.

\fait\ La parité avec M4.2 est vérifiée à $0{,}000\mathrm{e}{+}00$ d'écart
flottant (pas seulement en tolérance $10^{-9}$) sur deux régimes M4.2 et un
régime baseline M4.2B, quand \code{target\_rule="geometric"},
$\rho=1$, $\beta=1$ (\code{tests/test\_parity\_m4\_2.py}) : le refactor du
moteur (séparation taux/principal, table \code{loan\_events}) n'introduit
aucune dérive dynamique.

\subsection{Pathologie de service prédite et vérifiée}

\fait\ Pour un newborn $K_b=K_0=25$ face à une prêteuse
$K_\ell=\Kaut(0{,}5)=9801$ ($\gamma=0{,}5$, $A=1$, $\delta=0{,}01$) :
$r\approx0{,}02247$, $q_A\approx4888$, service $c=rq\approx109{,}85$,
production de l'emprunteuse au pas suivant
$F_{0,5}(K_b')\approx70{,}09$ : $c/F_{0,5}(K_b')\approx1{,}57>1$. Le
service excède la production du pas.

\fait\ Malgré cela, \textbf{0 défaut de liquidité n'a été observé sur
l'ensemble des runs pilotes} (\code{roots\_liquidity}=0 cumulé). La raison :
le paiement des intérêts est prélevé sur le \emph{stock} de capital $K$
(qui inclut le principal $q$ tout juste reçu, $\approx4913$), pas sur le
seul flux de production --- le stock suffit très largement. La fragilité
réelle est un problème de \textbf{levier} : la valeur nette de
l'emprunteuse ne change pas à la création du prêt ($K_b'-q_A=K_0=25$),
mais son ratio valeur-nette/dette tombe à $\approx0{,}5\,\%$, rendant
l'entité extrêmement vulnérable au \emph{prochain} choc $\xi$, pas au
service du prêt lui-même. C'est une fragilité d'insolvabilité différée,
pas une incapacité de service immédiate --- réponse à la question §21-Q10.

\section{Stationnarité : une horloge corrigée en cours de route}

\hyp\ initiale (avant tout run) : par analogie avec le temps de relaxation
autarcique individuel $K\mapsto(1-\delta)(K+AK^\gamma)$
($584$ pas pour atteindre $90\,\%$ de $\Kaut=9801$, $1048$ pas pour
$99\,\%$ à $\delta=0{,}01$), un horizon $T\approx8000$--$10000$ semblait
nécessaire.

\fait\ Cette hypothèse est \textbf{infirmée empiriquement}. Le run baseline
($T=3000$, seed 0) montre une population qui atteint la stationnarité dès
$t\approx375$--$750$ : moyenne de population $1134$--$1140$ sur les
demi-fenêtres $[750,1875]$ et $[1875,3000]$, coefficient de variation
$0{,}034$, dérive relative entre demi-fenêtres $<1\,\%$
(\code{lib\_metrics.window\_series\_metrics}). \inference : la mortalité
par cascades de faillites régule la population bien avant qu'aucune entité
individuelle n'approche $\Kaut$ --- le bon temps caractéristique est celui
de la dynamique de crédit-cascade, pas celui de la relaxation productive
autarcique. Conséquence pratique : $T=1000$ (burn-in $375$) suffit pour
l'exploration ; $T=3000$ est réservé aux cellules de confirmation.

\section{Objectif A : distribution des revenus d'intérêt}

\subsection{Corps et queue à la baseline}

\fait\ Sur 46 snapshots de la fenêtre stationnaire ($t\in[750,3000]$,
run baseline) : masse en zéro $p_0=0{,}089\pm0{,}008$ ; parmi les 13
familles de la ladder (\code{scripts/families.py}), \textbf{GB2 (4
paramètres) gagne systématiquement par AIC}, devant \code{lognorm\_mix\_5p}
et \code{gengamma\_3p} ; \code{pareto\_1p} (ajustée sur tout le support)
est très largement rejetée ($\Delta\mathrm{AIC}>2900$) ; le raccord
explicite lognormale+Pareto (\code{lognorm\_pareto\_mix\_5p}) est
\emph{systématiquement moins bon} que GB2 ou même le simple mélange de
lognormales, malgré le même nombre de paramètres (table
\ref{tab:aic}).

\begin{table}[h]
\centering
\caption{AIC des familles sur 3 snapshots représentatifs (intérêts reçus,
run baseline). Plus bas = meilleur.}
\label{tab:aic}
\begin{tabular}{lrrr}
\toprule
Famille (k paramètres) & $t=1500$ & $t=1550$ & $t=1600$ \\
\midrule
\code{expon\_1p} (1) & 9672,0 & 9031,8 & 9883,7 \\
\code{lognorm\_2p} (2) & 9973,9 & 9412,5 & 10207,5 \\
\code{gengamma\_3p} (3) & 9673,2 & 9030,0 & 9883,7 \\
\textbf{\code{gb2\_4p}} (4) & \textbf{9591,9} & \textbf{8917,7} & \textbf{9793,1} \\
\code{lognorm\_mix\_5p} (5) & 9610,8 & 8971,4 & 9809,8 \\
\code{lognorm\_pareto\_mix\_5p} (5) & 9829,3 & 9201,1 & 10015,7 \\
\code{pareto\_1p} (1) & 12544,8 & 12891,9 & 13538,2 \\
\bottomrule
\end{tabular}
\end{table}

\fait\ Le seuil de queue sélectionné par scan KS (\code{tail\_test}) donne
$\hat\alpha_{\mathrm{density}}=3{,}80\pm0{,}18$ (moyenne/écart-type sur les
46 snapshots) --- serré \emph{entre} snapshots, mais \textbf{instable au
seuil} : à seuil moitié, $\hat\alpha$ retombe systématiquement à
$2{,}52\pm0{,}12$. Un test de Vuong restreint à la queue,
loi-de-puissance-vs-exponentielle (\code{interest\_income.
exponential\_vs\_powerlaw\_lr}, famille explicitement demandée par le
cahier des charges), favorise la loi de puissance dans 46/46 snapshots
($z=3{,}15\pm0{,}87$) : la queue haute n'est pas un simple prolongement
exponentiel, mais son exposant n'est pas stable --- signature classique
d'une courbure log-log (corps à décroissance sous-exponentielle,
correctement capturée par GB2) plutôt que d'un régime Pareto authentique.
\inference : le critère 6 du §11 (stabilité au seuil) échoue ; il n'y a
\textbf{pas de régime Pareto robuste} à la baseline.

\subsection{Décomposition mécanique}

Avec $I_i=\sum_j q_{ij}r_{ij}$ sur le réseau contractuel, l'identité exacte
\[
\mathrm{Var}(\log I)=\mathrm{Var}(\log\,\mathrm{deg}_{\mathrm{out}})
+\mathrm{Var}(\log\,\overline{rq})
+2\,\mathrm{Cov}(\log\,\mathrm{deg}_{\mathrm{out}},\log\,\overline{rq})
\]
donne, moyennée sur la fenêtre stationnaire : \textbf{77\,\% de la variance
du log-revenu vient du nombre de prêts émis} (\code{deg\_out}), 23\,\% du
$rq$ moyen par prêt, covariance négligeable ($-0{,}2\,\%$). \fait\ Le taux
$r$ lui-même varie peu (moyenne $0{,}01797$, médiane $0{,}01752$, q90
$0{,}01865$ : bande de $6\,\%$) : le revenu est essentiellement
$I_i\approx \bar{r}\bar{q}\times\mathrm{deg}_{\mathrm{out}}(i)$.

\fait\ Corrélation de Spearman(revenu, âge)$\,\approx0{,}89$, très stable
entre snapshots (écart-type $0{,}006$) ; régression linéaire
revenu$\sim$âge : $R^2\approx0{,}72$. La distribution des âges rejette
formellement (test KS, $p\approx0$ dans les 46 snapshots) une loi
exactement exponentielle, mais reste sur-dispersée (écart-type/moyenne
$\approx2{,}1$, contre $1$ pour l'exponentielle) --- cohérent avec une
mortalité par cascades épisodiques plutôt qu'à taux constant.
\inference : le revenu d'intérêt est essentiellement un
\textbf{effet d'accumulation de contrats avec l'âge}, sous une mortalité
non exactement exponentielle mais sans excès extrême de dispersion --- pas
un phénomène de concentration par la taille des prêts individuels.

\fait\ Gini$(K)$ final $=0{,}054$ à la baseline : le capital est presque
parfaitement homogénéisé. C'est la matière première manquante pour un
$rq$ hétérogène.

\subsection{Le contrôle géométrique change l'interprétation}

Une première lecture (provisoire, avant ce run) attribuait le
\code{branching\_ratio} élevé ($0{,}786$) et le Gini$(K)$ bas ($0{,}054$) à
l'institution arithmétique elle-même. \fait\ Un run de contrôle, mêmes
paramètres sauf \code{target\_rule="geometric"}, T=1000 : branching
$=0{,}756$ (quasi identique à $0{,}787$), Gini$(K)=0{,}053$ (identique).
\textbf{Cette lecture initiale est infirmée} : la cause réelle est la
baseline lente $\delta{=}0{,}01,\sigma{=}0{,}01$ combinée à $\eta(N){=}N$,
pas la règle de principal.

En revanche, à cette même baseline, la règle géométrique donne
$\hat\alpha_{\mathrm{density}}=2{,}94$ contre $3{,}86$ en arithmétique
--- \textbf{une queue plus lourde}, avec une dominance du degré encore plus
forte ($0{,}872$ contre $\approx0{,}79$--$0{,}80$). Réponse à la question
§21-Q1 : le changement constitutif de M4.2B \textbf{rend l'objectif A plus
difficile}, pas plus facile, que M4.2.

\subsection{Balayage $\sigma$ : identification de mécanisme, pas contrôle
de queue}

\begin{table}[h]
\centering
\caption{Balayage $\sigma$ (arithmétique, $\rho=1$, $\lambda=30$, $T=1000$,
2 graines). Moyennes des deux graines.}
\label{tab:sigma}
\begin{tabular}{lrrrrrr}
\toprule
$\sigma$ & pop finale & Gini$(K)$ & part deg\_out & $\hat\alpha$(KS) & $\hat\alpha$(seuil/2) & $\hat\tau$ aval. \\
\midrule
0,01 & 1084 & 0,053 & 0,796 & 3,86 & 2,62 & 1,42 \\
0,03 & 927  & 0,065 & 0,775 & 4,15 & 2,72 & 1,69 \\
0,05 & 794  & 0,070 & 0,770 & 4,45 & 2,74 & 1,80 \\
0,10 & 630  & 0,092 & 0,742 & 5,03 & 2,71 & 2,00 \\
\bottomrule
\end{tabular}
\end{table}

\fait\ $\sigma$ agit dans le sens prédit sur les variables intermédiaires
(Gini$(K)$ et part deg\_out), mais $\hat\alpha$(KS) \textbf{s'alourdit}
avec $\sigma$ (table \ref{tab:sigma}) tandis que $\hat\alpha$(seuil/2) reste
quasi constant ($\approx2{,}6$--$2{,}7$). \inference : $\sigma$ réduit la
population stationnaire et l'espérance de vie effective, ce qui
\textbf{tronque l'accumulation de degré} --- le mécanisme même qui produit
la dispersion du revenu. Les deux effets ($\sigma$ injecte de
l'hétérogénéité de capital \emph{et} tronque l'accumulation de contrats)
s'opposent ; la troncature l'emporte dans la plage testée. Ce n'est pas un
échec du levier $\sigma$ mais une \textbf{identification de mécanisme}
conforme au critère 7 du §11 : un effet qui se réduit à un déplacement de
coupure n'est pas un contrôle de queue. Côté avalanches, $\sigma$ éloigne
le système d'un régime critique (branching $0{,}79\to0{,}46$, $\hat\tau$
$1{,}42\to2{,}00$) : pas de compromis favorable trouvé.

\section{Objectif B : avalanches de faillites}

\fait\ Test de Vuong tronquée-vs-lognormale sur le run baseline (fenêtre
stationnaire, \code{lib\_metrics.vuong\_trunc\_vs\_ln}, le modèle PRIMAIRE
--- pas la loi pure déjà rejetée par LRT) : $z=+16{,}3$, favorise très
fortement la loi de puissance tronquée. \code{branching\_ratio}
$=0{,}786$, stable entre demi-fenêtres.

\fait\ Balayage $\lambda\in\{10,100\}$ (2 graines, $T=1000$, sinon
baseline) : $\hat\tau$ dérive de $1{,}19$--$1{,}27$ ($\lambda=10$) à
$1{,}42$--$1{,}45$ ($\lambda=30$) à $1{,}75$ ($\lambda=100$), avec
\code{size\_max} $75\to155\to280$. \textbf{Vérification de l'artefact
avant conclusion} (cf. mémoire M4.2 sur ce piège précis) : la coupure
ajustée par \code{fit\_powerlaw\_cutoff} vaut $35$--$107$ (bien en-deçà de
\code{size\_max}) à $\lambda=10$ et $\lambda=30$, mais
\textbf{$\approx4{,}85\times10^{8}$ à $\lambda=100$ --- exactement
$e^{20}$, la borne numérique du paramètre \code{log\_cutoff} dans
\code{lib\_metrics.py}}. L'optimiseur a buté sur sa borne : aucune coupure
finie n'est détectée à $\lambda=100$, et le $\hat\tau=1{,}75$ rapporté est
en réalité l'exposant de la loi \textbf{pure}, non comparable aux deux
autres cellules (ajustements tronqués authentiques). \inference :
\textbf{aucun exposant indépendant de la taille du système n'est établi}
sur $10\le\lambda\le100$ avec cette méthode ; noter en outre que $\lambda$
déplace simultanément la taille du système et l'intensité de marché
($\eta(N)=N$), donc la dérive observée ne peut pas être attribuée à la
taille seule.

\fait\ Le \code{branching\_ratio} reste en revanche remarquablement stable
: $0{,}782$--$0{,}788$ sur $\lambda\in\{10,30,100\}$ \emph{et}
$0{,}756$--$0{,}787$ sous les deux \code{target\_rule}. \inference :
pinning institutionnel réel, analogue à celui trouvé par M4.2
($b\approx0{,}30$, \og un round de marché par tête\fg{}) mais à un niveau
différent ($\approx0{,}78$--$0{,}79$), déterminé par la baseline lente
$\delta{=}0{,}01,\sigma{=}0{,}01$ plutôt que par $\gamma$ ou la règle de
principal.

\section{Réponses au questionnaire du prompt (§21) disposant d'un support
empirique}

\begin{enumerate}[leftmargin=1.6em]
\item \textbf{Que change la cible arithmétique ?} Principal non borné plus
  élevé ($q_A/q_{\mathrm{geo}}=(1+\sqrt{K_\ell/K_b})/2$) ; égalisation
  exacte (auto-terminaison de la paire, \code{merge\_share}$\approx2{,}8\,\%$
  sous les deux règles) ; rend l'objectif A \emph{plus} difficile.
\item \textbf{Forme du corps ?} GB2, systématiquement meilleur par AIC.
\item \textbf{Queue Pareto robuste ?} Non, dans aucune configuration
  testée.
\item \textbf{Famille globale ?} GB2 ; le composite Pareto explicite
  n'apporte rien.
\item--\item[8] \emph{(domaine des exposants, rôle de $\eta$)} :
  \textbf{non tranchés} --- $\rho$ testé seulement sur $\{0{,}125,1,2\}$,
  une graine, $T$ court ; aucune tendance monotone claire.
\item \textbf{Rôle de $r,q,rq$ ?} $r$ quasi constant ; 77\,\% de la
  variance du log-revenu vient du nombre de contrats, 23\,\% du $rq$ moyen.
\item \textbf{Pathologie de la règle de taux ?} Oui en levier (ratio
  valeur-nette/dette $\approx0{,}5\,\%$ pour un newborn face à une
  entité mature), non en service (0 défaut de liquidité observé).
\item \textbf{Avalanches en loi de puissance ?} Branching robuste et
  pinné ; exposant indépendant de la taille non établi.
\item \textbf{Compromis queue/criticalité ?} $\sigma$ dégrade les deux
  simultanément dans la plage testée --- pas de compromis favorable trouvé.
\item[13] \emph{(domaine conjoint)} : \textbf{non atteint} dans la plage
  testée.
\item[15] \textbf{Coupure, échelle ou exposant ?} Coupure, de façon
  répétée et démontrée ($\sigma$, $\lambda$ déplacent tous une coupure
  haute, jamais la forme du corps).
\item[16] \emph{(refonte nécessaire ?)} Aucune refonte de mécanisme n'a
  été tentée dans cette phase ; le levier testé qui identifie le mécanisme
  le plus clairement est $\sigma$ (troncature de l'accumulation de degré),
  pas une institution à remettre en cause.
\end{enumerate}

\section{Non testé (§18)}

$\gamma\ne0{,}5$ ; $K_0\ne25$ ; $\eta$ non linéaire ($\beta\ne1$) ;
ablations de mécanisme (règle de taux alternative, fusion/mémoire des
contrats) ; campagne multi-graines complète par cellule ; collapse de
taille finie en bonne et due forme pour les avalanches (nécessiterait de
découpler taille du système et intensité de marché, et un $\lambda>100$
avec une coupure correctement dans la plage observée). Le budget de calcul
de cette phase a été alloué aux pilotes ci-dessus sur la base du
diagnostic qu'ils apportaient l'information marginale la plus grande avant
de figer un protocole de campagne.

\end{document}
